\section{Uniform phase separation and positive-measure coercivity}
\label{sec:periodic-coercivity}
We strengthen the frequency estimate by comparing consecutive exact
minimizers. The choice of a period will depend on the frequency and
possibly on the radius. This dependence is essential: a fixed triple
of periods is not asserted to separate the entire unbounded frequency
axis by a positive constant.

Retain $K_R(y,z)=\ell_R(y,z)-g$ and $W_{m,R}$ from
\eqref{eq:half-action-v4}. Write
\[
 a_m(R)=-y_m^{(m,R)}>0,\qquad
 \Delta_m(R)=E_{m+1}(R)-E_m(R)>0.
\]
The symbol $a_m$ denotes a terminal height; it is unrelated to the
logarithmic constant in \eqref{eq:frequency-period-budget}.

\subsection{Comparison of successive exact minima}
\begin{lemma}[The terminal value function]
\label{lem:terminal-response}
For $y\in[-d,0]$ put
\[
 V_R(y)=\min_{z\in[-d,0]}\{K_R(y,z)+u_R(z)\},\qquad
 \delta_R(y)=V_R(y)-u_R(y).
\]
The minimizer $f_R(y)$ is unique. It belongs to $(-d,0)$ for $y<0$
and equals zero for $y=0$. The functions are twice continuously
differentiable, with one-sided derivatives at the endpoints, and
\begin{equation}\label{eq:terminal-response}
 f_R'(y)\ge\frac1{14},\qquad
 |f_R(y)|\ge\frac{|y|}{14},\qquad
 |\delta_R''(y)|\le2,\qquad |\delta_R'(y)|\le2|y|.
\end{equation}
These estimates hold on the entire height interval, uniformly in $R$.
\end{lemma}
\begin{proof}
The second derivative in $z$ is $K_{zz}+u''>0$. At $z=-d$,
\[
 K_z(y,-d)=\frac{h(y,-d)(-d)/x_R(-d)-d-y}{\ell_R(y,-d)}<0,
 \qquad u_R'(-d)<0.
\]
At $z=0$ and $y<0$ the derivative is $-y/\ell_R(y,0)>0$.
Consequently the minimizer is interior and unique. For $y=0$,
\eqref{eq:kernel-coercivity} gives the unique minimizer zero.
The implicit function theorem applies to the stationarity equation,
including its smooth extension near the endpoints, and gives
\[
 f_R'(y)=\frac{-K_{yz}(y,f_R(y))}
                    {K_{zz}(y,f_R(y))+u_R''(f_R(y))}
 \ge\frac{1/2}{4+3}=\frac1{14}.
\]
Here the individual segment bounds established in the proof of
Lemma~\ref{lem:uniform-half-hessian} were used. Integration from $y$
to zero proves the second inequality of \eqref{eq:terminal-response}.

The envelope formula is
\[
 V_R''(y)=K_{yy}(y,f_R(y))-
       \frac{K_{yz}(y,f_R(y))^2}{K_{zz}(y,f_R(y))+u_R''(f_R(y))}.
\]
The Hessian of $K_R(y,z)+u_R(z)$ dominates $(3/2)I_2$.
Its Schur complement therefore is at least $3/2$: evaluate its
quadratic form on $(1,t)$ and minimize over $t$. Also
$V_R''\le K_{yy}\le4$. Since $100/47\le u_R''<3$, subtraction
shows $|\delta_R''|\le2$. The first derivatives of $K_R$ and $u_R$
vanish at the origin, so $\delta_R'(0)=0$. Integration proves the
last inequality. All estimates concern the nonlinear action itself.
\end{proof}

\begin{proposition}[Consecutive terminal heights]
\label{prop:terminal-comparison}
For every $m\ge1$ and $R\in I$,
\begin{equation}\label{eq:terminal-comparison}
 a_{m+1}(R)\ge\frac3{70}a_m(R).
\end{equation}
\end{proposition}
\begin{proof}
Let $y$ minimize $W_{m,R}$ and let $(X,z)$ minimize $W_{m+1,R}$,
where $X$ consists of the first $m$ coordinates. Eliminating the
last coordinate gives exactly
\[
 \min_z W_{m+1,R}(X,z)=W_{m,R}(X)+\delta_R(X_m),
 \qquad z=f_R(X_m).
\]
All minimizing coordinates are interior by
Lemma~\ref{lem:excursion-family}. Stationarity and the fundamental
theorem of calculus thus yield
\[
 A(X-y)=-\delta_R'(X_m)e_m,\qquad
 A=\int_0^1\nabla^2W_{m,R}(y+t(X-y))\dd t\ge3I_m.
\]
The segment stays in the height box. In particular
$0<(A^{-1})_{mm}\le1/3$. Taking the last coordinate and applying
Lemma~\ref{lem:terminal-response}, we obtain
\[
 |y_m|\le |X_m|+\frac13|\delta_R'(X_m)|
             \le\frac53|X_m|.
\]
Finally $|z|=|f_R(X_m)|\ge|X_m|/14$. These two inequalities give
\eqref{eq:terminal-comparison}. No monotonicity of the terminal
heights, and no quadratic replacement of $W_{m,R}$, is assumed.
\end{proof}

\begin{corollary}[A lower ratio for the excess increments]
\label{cor:increment-ratio}
For every $m\ge1$ and $R\in I$,
\begin{equation}\label{eq:increment-ratio}
 \Delta_{m+1}(R)\ge10^{-5}\Delta_m(R).
\end{equation}
\end{corollary}
\begin{proof}
The append and truncate comparisons in the proof of
Theorem~\ref{thm:quantitative-excess} give
\[
 \Delta_m\le\frac{a_m^2}{g},\qquad
 \Delta_{m+1}\ge\frac2R a_{m+2}^2.
\]
Applying \eqref{eq:terminal-comparison} twice, and using
$g>79/100$ and $R\le47/100$, yields
\[
 \Delta_{m+1}\ge\frac{2g}{R}\left(\frac3{70}\right)^4\Delta_m
 >\frac{158}{47}\left(\frac3{70}\right)^4\Delta_m
 >\frac{\Delta_m}{100000}.
\]
This is a comparison between adjacent increments, stronger for the
present purpose than separate exponential upper and lower bounds.
\end{proof}

\subsection{A frequency-dependent triple and a fixed finite menu}
Put
\begin{equation}\label{eq:uniform-phase-constants}
 \eta_0=\frac1{4000000},\qquad
 \sigma_0=\frac1{6000000\pi},\qquad
 M(b)=m(b)+1,
\end{equation}
where $m(b)$ is defined in \eqref{eq:frequency-period-budget}.

\begin{theorem}[Uniform positive phase separation]
\label{thm:uniform-phase-separation}
For every $R\in I$ and $b\ge1$ the integer
\begin{equation}\label{eq:adaptive-period}
 j(R,b)=\min\{j\ge1:b\Delta_j(R)\le1/2\}
\end{equation}
is well defined and satisfies $j(R,b)\le m(b)$. If $|\beta|=b$,
then, for arbitrary $u,s,c$,
\begin{equation}\label{eq:uniform-phase-separation}
 \max_{P\in\{P_0,P_{j(R,b)},P_{j(R,b)+1}\}}
       |Z_P(u,s,\beta,c;R)-1|\ge\sigma_0.
\end{equation}
The longest of these physical words has at most $2m(b)+3$ collisions.
For the entire band $1\le|\beta|\le B$, the fixed menu
\begin{equation}\label{eq:fixed-period-menu}
 \mathcal P_B(R)=\{P_0(R),P_1(R),\ldots,P_{m(B)+1}(R)\}
\end{equation}
satisfies the same lower bound for its maximum, uniformly in $R$.
It contains $m(B)+2$ periods.
\end{theorem}
\begin{proof}
The upper bound of \eqref{eq:quantitative-excess} gives
$b\Delta_{m(b)}\le1/300$. Thus the defining set in
\eqref{eq:adaptive-period} is nonempty. If $j(R,b)=1$, its lower
bound gives $b\Delta_1\ge\Delta_1\ge1/4000000=\eta_0$.
If $j(R,b)>1$, minimality and \eqref{eq:increment-ratio} give
\[
 b\Delta_{j(R,b)}>10^{-5}/2\ge\eta_0.
\]
In either case
\begin{equation}\label{eq:selected-phase-window}
 \eta_0\le b\Delta_{j(R,b)}\le1/2.
\end{equation}
The exact three-period identity \eqref{eq:three-period-cancellation}
therefore has distance at least $2\eta_0/\pi$ from one. The distance
of a product of three unit complex numbers from one is at most the
sum of their three distances. This proves
\eqref{eq:uniform-phase-separation}, since
$2\eta_0/(3\pi)=\sigma_0$. The physical and induced counts are
those of Lemma~\ref{lem:excursion-family}; hence the stated word
length. Finally $m(b)\le m(B)$ on the band, so the selected triple
is contained in \eqref{eq:fixed-period-menu}.

The index selection uses a first crossing, not monotonicity of
$\Delta_j$. It need not be continuous in $R$ or $b$. The menu has
fixed indices and consists of analytically varying true periods,
which is the asserted parameter-uniform conclusion.
\end{proof}

\begin{corollary}[Logarithmic uniform-norm coercivity]
\label{cor:logarithmic-coercivity}
Let $q$ be circle-valued and measurable on $Y_R^*$, and continuous
in neighborhoods of every state of the triple in
\eqref{eq:uniform-phase-separation}. With $D(q)$ as in
Corollary~\ref{cor:approximate-phase-bound}, for $b=|\beta|\ge1$,
\begin{equation}\label{eq:logarithmic-coercivity}
 D(q)\ge\frac1{4000000\pi(j(R,b)+1)}
       \ge\frac1{4000000\pi M(b)}.
\end{equation}
\end{corollary}
\begin{proof}
Continuity and full support extend the essential bound to the
selected regular periodic states. Telescoping around each period
bounds its phase discrepancy by its induced length times $D(q)$.
The three induced lengths sum to $2(j(R,b)+1)$. Apply the product
identity and \eqref{eq:selected-phase-window} to obtain
\[
 2\eta_0/\pi\le2(j(R,b)+1)D(q).
\]
This proves \eqref{eq:logarithmic-coercivity}. The modulus-one
condition is used in telescoping and is not a consequence of an
arbitrary operator normalization.
\end{proof}

\subsection{Regular tubes of the actual return map}
We use the Euclidean metric in local $(\alpha,p)$ coordinates. The
following uniformity concerns the specified periodic family, not all
branches of the induced map.

\begin{lemma}[Uniform return neighborhoods of the period family]
\label{lem:uniform-period-tubes}
There are constants $0<\rho_0<1$, $L\ge2$ and $B_0\ge1$, independent
of $R$ and $m$, with the following properties. For every selected
state $x_P$ of $P_0$ or any $P_m$, the ball $B(x_P,\rho_0)$ lies
in $Y_R^*$ and in one regular first-return branch. On this ball
$r_R^*$ and $\kappa_R^*$ are constant, $r_R^*\in\{2,3\}$, and
\[
 \operatorname{Lip}(F_R^*)\le L,\qquad
 \operatorname{Lip}(\tau_R^*)\le B_0.
\]
If $P$ has induced period $h$ and $0<\rho\le\rho_0L^{-h}$, every
$x\in B(x_P,\rho)$ follows the same $h$ return branches as $x_P$,
and, writing $F=F_R^*$,
\begin{align}\label{eq:period-tube-bounds}
 d(F^k x,F^k x_P)&\le L^k\rho &&(0\le k\le h),\\
 |T_{h,R}(x)-T_{h,R}(x_P)|&\le2B_0L^h\rho,\\
 \nu_R^*(B(x_P,\rho))&=\rho^2/(4c_*).
\end{align}
\end{lemma}
\begin{proof}
We verify that the family stays a uniform distance from every local
singularity relevant to a single return. The estimates in
Lemma~\ref{lem:excursion-family} hold for all its contact heights:
$-2d/5<y<0$, the long-flight third-disk clearances are at least
$9/500$, and the incidence cosine at $A$ exceeds $1/25$.
At the facing contacts on $O,C$ it is bounded below by a positive
constant from the horizontal-component inequalities in that proof.
End flights have a uniform positive distance from every third disk:
the opposite facing center is horizontally separated by $D/2$,
$B$ is vertically separated, and the remaining lattice centers obey
the enumerated uniform separation bounds. All flight lengths are
bounded above and below by positive constants.

Return membership has uniform margins as well. The selected contacts
on $O$ satisfy
\[
 |\alpha-\pi/6|\le1/22<3/50,\qquad
 |p|\le1/22+2/79<3/20.
\]
The unselected contacts on $C$ and $A$ have angles near $7\pi/6$
and equal to $2\pi/3$, respectively, uniformly separated from $U$
and the four squares in $Y_R$. The long normal orbit has the same
strict margins. Each selected-to-selected path consequently contains
two physical flights, or three through $C,A,O$.

For clarity, compactness is applied only to these bounded-length
paths. The parameters $R$ and their contact coordinates range in a
compact box. Every limiting path satisfies the non-grazing,
nonocclusion and return-membership margins just established. In a
neighborhood of such a path the selected entry roots are simple;
the physical collision map and flight time are analytic in the
initial state and parameter. There are only the finitely many
center itineraries just listed, up to lattice translation. A finite
cover of their compact closures therefore gives a common neighborhood
radius and uniform first-derivative bounds for the actual two- or
three-flight first returns. Decrease the radius so its coordinate
balls lie in $Y_R^*$ and contain no change of lattice labels; enlarge
the derivative bounds to $L\ge2$, $B_0\ge1$. This proves the first
part without taking a derivative of an arbitrarily long return.

Starting in a ball of radius at most $\rho_0L^{-h}$, induction gives
the first inequality in \eqref{eq:period-tube-bounds} and preserves
the same return decisions at each step. Summing the roof variations
gives $B_0\rho\sum_{k<h}L^k\le2B_0L^h\rho$. Finally
$\dd\nu_R^*=(4\pi c_*)^{-1}\dd\alpha\dd p$ on the ball, whose
Euclidean area is $\pi\rho^2$. This proves the last equality.
\end{proof}

\begin{theorem}[Positive-measure coercivity for regular phases]
\label{thm:positive-measure-coercivity}
Fix $0<\alpha\le1$, $H\ge1$ and $1\le p<\infty$. Suppose $q$ is
measurable and circle-valued on $Y_R^*$ and has an $\alpha$-H\"older
representative, of constant at most $H$, on the balls of radius
$\rho_0$ around all states of the selected triple. Define
\[
 \mathcal D_q(x)=e^{i(u\cdot\kappa_R^*(x)+s r_R^*(x)+\beta\tau_R^*(x))}
                      q(F_R^*x)-e^{ic}q(x).
\]
For $b=|\beta|\ge1$, put
\begin{equation}\label{eq:coercivity-radius}
 \rho(b,H)=L^{-M(b)}
 \min\left\{\frac{\rho_0}{2},\frac{\sigma_0}{8bB_0},
                       \left(\frac{\sigma_0}{8H}\right)^{1/\alpha}\right\}.
\end{equation}
Then the actual invariant return probability satisfies
\begin{equation}\label{eq:positive-measure-coercivity}
 \|\mathcal D_q\|_{L^p(\nu_R^*)}
 \ge\frac{\sigma_0}{2M(b)}
             \left(\frac{\rho(b,H)^2}{4c_*}\right)^{1/p}.
\end{equation}
All geometric constants are uniform over $R\in I$. No independence
of the returns is assumed.
\end{theorem}
\begin{proof}
Select from the triple a period $P$ with $|Z_P-1|\ge\sigma_0$.
Let its selected state be $x_P$ and its induced period be $h\le M(b)$.
Write $\varphi=u\cdot\kappa_R^*+s r_R^*+\beta\tau_R^*$, $F=F_R^*$
and $S_h\varphi=\sum_{k<h}\varphi\circ F^k$. The lattice and count
terms are constant along the corresponding branches in the tube.
Thus for $x\in B(x_P,\rho(b,H))$, Lemma~\ref{lem:uniform-period-tubes}
gives
\[
 |S_h\varphi(x)-S_h\varphi(x_P)|\le2bB_0L^h\rho(b,H).
\]
At $x_P$, the modulus of
$e^{iS_h\varphi(x_P)}q(F^h x_P)-e^{ich}q(x_P)$ is $|Z_P-1|$.
Since $F^h x_P=x_P$, the H\"older bounds at the two endpoints and
\eqref{eq:coercivity-radius} imply, throughout the smaller ball,
\begin{align*}
 \left|e^{iS_h\varphi(x)}q(F^h x)-e^{ich}q(x)\right|
 &\ge\sigma_0-2bB_0L^h\rho(b,H)-2H(L^h\rho(b,H))^\alpha\\
 &\ge\sigma_0/2.
\end{align*}
Choose the indicated representative on these balls; changes on null
sets do not affect the following integrals. Telescoping $h$ times
bounds the left side by $\sum_{k<h}|\mathcal D_q(F^k x)|$ almost
everywhere. Minkowski's inequality on the smaller ball and invariance
of $\nu_R^*$ give
\[
 \frac{\sigma_0}{2}\nu_R^*(B(x_P,\rho(b,H)))^{1/p}
 \le\sum_{k<h}\|\mathcal D_q\circ F^k\|_{L^p(B(x_P,\rho(b,H)),\nu_R^*)}
 \le h\|\mathcal D_q\|_{L^p(\nu_R^*)}.
\]
Use the exact ball mass and $h\le M(b)$ to conclude. The estimate
uses measure preservation, not disjointness or probabilistic
independence of the iterates of the ball.
\end{proof}

\begin{corollary}[A polynomial budget for the phase regularity]
\label{cor:polynomial-regularity-budget}
Fix $\alpha,p,H_0$ as above and $\zeta\ge0$. If the H\"older constants
satisfy $H\le H_0b^\zeta$, set
\begin{equation}\label{eq:coercivity-exponent}
 \kappa=\frac{\log L}{\log(100/49)}+
                       \max\{1,\zeta/\alpha\}.
\end{equation}
There is $C>0$, depending only on these fixed constants and the
mechanical family, such that
\begin{equation}\label{eq:polynomial-coercivity}
 \|\mathcal D_q\|_{L^p(\nu_R^*)}
             \ge C\frac{b^{-2\kappa/p}}{1+\log b}\qquad(b\ge1).
\end{equation}
\end{corollary}
\begin{proof}
Writing $a=\log(100/49)$, we have $M(b)\le3+\log b/a$ and hence
$L^{-M(b)}\ge L^{-3}b^{-\log L/a}$. The minimum in
\eqref{eq:coercivity-radius} is at least
\[
 \min\left\{\frac{\rho_0}{2},\frac{\sigma_0}{8B_0},
           \left(\frac{\sigma_0}{8H_0}\right)^{1/\alpha}\right\}
                     b^{-\max\{1,\zeta/\alpha\}}.
\]
Substitution into \eqref{eq:positive-measure-coercivity} proves the
claim with a positive constant independent of $R$ and $b$.
\end{proof}

\subsection{The operator interface}
Theorems~\ref{thm:uniform-phase-separation} and
\ref{thm:positive-measure-coercivity} have different hypotheses.
The former is an unconditional estimate for the specified mechanical
periods. The latter is an estimate for a stated regularity class of
circle-valued phases, in the actual invariant measure. In particular
it supplies a positive-measure obstruction rather than a conclusion
obtained by evaluating an unspecified measurable representative at
isolated periodic points.

To apply \eqref{eq:polynomial-coercivity} to a transfer operator,
one must construct from its approximate spectral vectors phases in
that class, bound the H\"older constant, and compare their physical
$L^p$ defect with the operator defect. Such a comparison would
contradict an error smaller than the right side of
\eqref{eq:polynomial-coercivity}. It is not proved merely by writing
an eigenvalue equation. In particular a division by the modulus of a
spectral vector requires a separate lower-modulus estimate.
General Liv\v{s}ic regularity results for Markov systems
\cite{BHN} concern precise dynamical and cohomological hypotheses;
they do not supply a uniform quantitative regularity bound for these
approximate phases without checking those hypotheses and the defect
comparison. No such comparison is assumed here for the full induced
operator. The common operator realization and the all-branch raw
residual estimates in the final section remain the requirements for
the original joint density theorem.
